\scanpage{24}
\uline{Espaces essentiels et ensembles ordonnés}

$X$ un ensemble

a) $\unsure{\tau|\mathcal T}(X)$ l'ens des topologies sur $X$

$\underline\omega(X)$ \rule{2.6em}{0.35pt} préordres sur $X$

\[
\deleted{\unsure{\alpha_X}}\quad \tau(X)\xrightarrow{\alpha_X}\underline\omega(X)
\]

à $T\in\tau(X)$ est associée la relation $\alpha_X(T)$

\[
y\leq_T x\mathrel{\mathop{\Longleftrightarrow}^{(\text{déf})}}x\in\overline{\{y\}}
\]

\[
\deleted{\unsure{\beta_X}}\quad\tau(X)\xleftarrow{\beta_X}\underline\omega(X)
\]

à la relation $\omega$ $({\leq}_{\omega})$ est associée $T$

\[
\begin{aligned}
&\deleted{\unsure{\{U\subset X\}|\mathrm{Fer}}}\\[-2pt]
\mathrm{Ouv}(T)&=\{\unsure{U|X}\subset X\mid(x\in U,\,y\leq_{\omega}x)\Rightarrow y\in U\}\\
&=\text{réunions quelconques d'ensembles}\\
&\phantom{=}\text{de la forme }\ O_x=(\unsure{]\leftarrow,x]})=\{y\in X\mid y\leq_{\omega} x\}
\end{aligned}
\]

\[
\mathrm{Fer}(T)=\{Z\subset X\mid(x\in Z,\,y\unsure{\geq|\leq}_{\omega}x)\Rightarrow y\in Z\}
\]

$\alpha$ et $\beta$ fonctoriels en l'ensemble $X$
(pour images inverses de top. resp. de relations d'ordre)

b) Soit $f:X\longrightarrow X'$ une application, \deleted{$X$ muni de $(T,\omega)$, $X'$ de $( $}
si $X,X'$ sont des espaces top. et si $f$ est \deleted{croissante}\inserted{continue}, $f$ est croissante
si $X,X'$ sont préordonnés et si $f$ est croissante, $f$ est continue

$\unsure{\therefore|\text{c)}}$ Deux foncteurs

\diagram{\begin{tikzcd}[ampersand replacement=\&]
\deleted{\unsure{\alpha}}\quad(\mathrm{Esp}) \arrow[r,shift right,"\alpha"'] \& (\mathrm{Préor}) \arrow[l,shift right,"\beta"']
\end{tikzcd}}

c) \deleted{\unsure{Définition de}}
Soit $T\in\tau(X)$, montrer que \unsure{pour} tout ouvert de $X$ est stable par générisation
(\deleted{\unsure{ou|resp.}} tout fermé de $X$ stable par spécialisation), i.e. ouvert pour $\dot\alpha_X(T)$, i.e.

\[
T\leq\beta_X\alpha_X(T)\deleted{\unsure{\leq}}
\]

i.e on a une application continue

\[
\beta\alpha(X)\longrightarrow X
\]

dont l'application sous-jacente est l'identité, d'où une transfo fonctorielle en $X$

\[
\beta\alpha\xrightarrow{\rho}\mathrm{id}_{(\mathrm{Esp})}
\]

Soit $\omega\in\underline\omega(X)$, \deleted{\unsure{munissons $X$ de}} $T=\beta_X(\omega)$, alors $x\in\overline{\{y\}}$ sss
le plus petit ouvert $O_x$ contenant $x$ contient $y$, i.e. sss $y\geq x$, donc on trouve

\[
\alpha_X\beta_X(\omega)=\omega
\]

i.e. on a un isomorphisme

\[
X\xrightarrow{\overset{\sigma}{\sim}}\alpha\beta(X)
\]

tel que l'application ensembliste sous-jacente soit l'identité, d'où

\[
\mathrm{id}_{(\mathrm{Préor})}\xrightarrow{\overset{\sigma}{\sim}}\alpha\beta
\]

d) Montrer que $\rho$ et $\sigma$ sont des morphismes d'adjonction, associés,
établissant pour

\[
X\in\mathrm{Ob}(\mathrm{Esp})\qquad I\in\mathrm{Ob}(\mathrm{Préor})
\]
% archive sheet number: 23
